\documentclass[11pt]{article}
\usepackage[margin=1in]{geometry}
\usepackage{amsmath,amssymb,amsthm,array}
\usepackage[colorlinks=true,linkcolor=blue,citecolor=blue,urlcolor=blue]{hyperref}
\newtheorem{theorem}{Theorem}[section]
\newtheorem{lemma}[theorem]{Lemma}
\newtheorem{proposition}[theorem]{Proposition}
\newtheorem{corollary}[theorem]{Corollary}
\newtheorem{remark}[theorem]{Remark}
\numberwithin{equation}{section}
\newcommand{\R}{\mathbb R}
\newcommand{\rank}{\operatorname{rank}}
\newcommand{\diag}{\operatorname{diag}}
\newcommand{\op}{\mathrm{op}}
\newcommand{\Frob}{\mathrm F}
\newcommand{\range}{\operatorname{range}}
\newcommand{\tr}{\operatorname{tr}}
\newcommand{\cB}{\mathcal B}
\newcommand{\cR}{\mathcal R}
\title{Bounded-spectrum fractional updates\protect\\
and relative principal-root corrections}
\author{Working research manuscript}
\date{November 2026}
\begin{document}
\maketitle

\begin{abstract}
We give two refinements of low-rank fractional matrix-function updates
from raw factors. If positive semidefinite endpoints have norms at most
$R$ times their perturbation norm, fixed $R>1$, a conventional
$\widetilde O_\alpha(Nmr+m^3)$ construction attains the sharp leading
width $r\log^2(1/\varepsilon)/(4\pi^2\alpha)$ for their fractional-power
correction. All large-shift quadrature terms share one polynomial Krylov
block; no best-rational-approximant oracle is used. For square roots we
also construct a correction with error relative to the correction itself,
including singular endpoints and signed updates. Its two-factor width
is $O(rL[L+\log\rho])$, where $L=\log(C/\varepsilon)$ and $\rho$ is
an explicitly computable endpoint-to-update scale parameter. Classical
low-rank ADI is executed through Stieltjes integrals over raw shifted
solves, with all action errors and preprocessing charged. A diagonal-Gram
promise removes cubic preprocessing deterministically. These are
exact-arithmetic results; the general raw-factor backend remains cubic
in the square case.
\end{abstract}

\noindent\textit{Review and priority status.}
The complete research note and typeset source passed separate independent
internal reviews. No external or formal verification, finite-word stability,
or certified novelty is claimed.
ADI, Sylvester displacement, Stieltjes integrals, polynomial Krylov
compression and randomized preconditioning are established methods.

\section{Statements and arithmetic model}

Let $X\in\R^{N\times m}$, $U\in\R^{N\times r}$ and $K=K^T\in
\R^{r\times r}$, with $1\le r\le m\le N$. Throughout,
\begin{equation}\label{eq:input}
 A=XX^T,\qquad E=UKU^T,\qquad B=A+E\succeq0.
\end{equation}
The norm without subscript is the operator norm. We count exact real
field operations, comparisons, positive scalar square roots and discrete
randomness. Complex or large matrix-function oracles are not inputs.
The notation $\widetilde O$ suppresses only logarithms of dimensions,
$1/\varepsilon$ and $1/\delta$. Additional scale logarithms will remain
explicit through $H$ below.
For general fixed $0<\alpha<1$, input-dependent scalar powers must be
added to the model or charged separately, as in the companion algorithm
\cite{Raw}. The square-root case needs only the operations just listed.

Thin QR gives $U=Q_0R_0$, with $Q_0$ orthonormal on its column support.
Define
\begin{equation}\label{eq:intrinsic}
 C=R_0KR_0^T,\qquad E=Q_0CQ_0^T,\qquad
 a=\|C\|_{\Frob}=\|E\|_{\Frob}.
\end{equation}
This costs $O(Nr^2+r^3)\le O(Nmr)$, and
\begin{equation}\label{eq:anorm}
 \|E\|\le a\le\sqrt r\|E\|.
\end{equation}
If $a=0$, every correction below is zero and the algorithm returns zero.
Henceforth $a>0$. Singular and arbitrarily cancelling supplied update
factors are allowed. Actual support sizes smaller than $r$ only improve
the stated bounds.

\begin{theorem}[Constructive sharp width at bounded endpoint scale]
\label{thm:bounded}
Fix $0<\alpha<1$, a positive integer $r$, and $R>1$. Under
\eqref{eq:input}, additionally promise
\begin{equation}\label{eq:Rpromise}
 \max(\|A\|,\|B\|)\le R\|E\|.
\end{equation}
For $\varepsilon,\delta\in(0,1)$ an algorithm returns a real signed
factor $F=Z\Sigma Z^T$, $\Sigma$ diagonal with entries $\pm1$, such
that, with probability at least $1-\delta$,
\[
 \|F-(B^\alpha-A^\alpha)\|\le\varepsilon\|E\|^\alpha.
\]
Its conventional work is
$(Nmr+m^3)$ times a fixed polynomial in $\log(2r/\varepsilon)$ and
the allowed dimension/confidence logarithms. Uniformly in finite
dimensions and all promised inputs, its width is at most
\begin{equation}\label{eq:sharpwidth}
 \left[\frac{r}{4\pi^2\alpha}+o_{\alpha,R,r}(1)\right]
                                      \log^2(1/\varepsilon).
\end{equation}
This leading coefficient is optimal by the bounded-spectrum minimax
theorem \cite{Bounded}. If $E\succeq0$, $F$ may be positive
semidefinite; if $E\preceq0$, it may be negative semidefinite.
Dense expansion costs an additional $O(N^2\,\mathrm{width})$.
\end{theorem}

The fixed promise $R>1$ is material. The class with $R=1$ has a
saturation exception and leading minimax coefficient
$(r-1)/(4\pi^2\alpha)$ \cite{Bounded}. We do not claim to attain that
boundary improvement here. For square roots, \eqref{eq:sharpwidth}
has coefficient $r/(2\pi^2)$, half the unrestricted-endpoint coefficient.

For the relative theorem, let $\Lambda>0$ be any known common endpoint
norm bound. For example, the raw input supplies the valid choice
$\Lambda=\|X\|_{\Frob}^2+a$ at cost $O(Nm)$. Put
\begin{equation}\label{eq:relative-scales}
 d_0=\frac{a}{2\sqrt{r\Lambda}},\qquad
 \rho=\frac{a}{d_0^2}=\frac{4r\Lambda}{a},\qquad
 L=\log(C_*/\varepsilon),
\end{equation}
where $C_*>1$ is an absolute constant.

\begin{theorem}[Relative principal-root correction]
\label{thm:relative}
For $\varepsilon,\delta\in(0,1)$ and \eqref{eq:input}, set
$\Delta=B^{1/2}-A^{1/2}$ and use \eqref{eq:relative-scales}.
An algorithm returns real two-factor output
$F=P_fQ_f^T$ with at most $kr$ columns in each factor, satisfying
\begin{equation}\label{eq:relative-error}
 \Pr\{\|F-\Delta\|\le\varepsilon\|\Delta\|\}\ge1-\delta,
 \qquad k=O\bigl(L[L+\log\rho]\bigr).
\end{equation}
Its symmetrization has the same error guarantee and at most $2kr$
signed columns. Let $H$ be a sufficiently large absolute multiple of
$L+\log(2\rho)+\log(2r)+1$. A conventional operation bound is
\begin{equation}\label{eq:relative-work}
 \widetilde O\bigl(NmrH^5+m^3H^4\bigr).
\end{equation}
Neither endpoint need be nonsingular. There is no smallest positive
eigenvalue or downdate-margin parameter. The scale dependence through
$\rho$ is explicit. Dense expansion adds $O(N^2kr)$.
\end{theorem}

The relative ADI construction does not assert PSD output, even if
$E\succeq0$, or optimal joint dependence on relative accuracy and
$\rho$. Its smaller width arises because internal quadrature stages
compute actions without creating extra final columns.

\begin{theorem}[Deterministic backend without cubic setup]
\label{thm:diagonal}
Let $D=\diag(\|X_{:1}\|^2,\ldots,\|X_{:m}\|^2)$ and promise
\begin{equation}\label{eq:diagonal-promise}
 \tfrac34D\preceq X^TX\preceq\tfrac54D.
\end{equation}
Then Theorems~\ref{thm:bounded} and \ref{thm:relative} hold
deterministically without cubic preprocessing. The relative algorithm
uses $O(NmrH^5)$ operations; the bounded-spectrum algorithm uses
$Nmr$ times a fixed polynomial in its accuracy logarithm.
Zero columns cause no problem. The promise permits arbitrarily unequal
column magnitudes and arbitrarily large nonzero spectral condition
numbers. It is not supplied or verified for free for arbitrary $X$.
\end{theorem}

For $N=m$, fixed $r$ and polynomially bounded $\rho$, the last theorem
gives ordinary $N^2$ arithmetic up to the stated logarithms, including
explicit output. Its input promise is essential. For arbitrary raw
factors, \eqref{eq:relative-work} remains cubic in the square case;
none of these statements establishes general quadratic multiplication.

\section{Common inverse operators and positive quadrature}

We state exactly the inputs taken from the companion algorithm
\cite{Raw}, and record the core algebra needed for both refinements.
For $t>0$ write
\[
 R_A(t)=(A+tI)^{-1},\quad R_B(t)=(B+tI)^{-1},\quad
 V_t=R_A(t)Q_0,\quad M_t=Q_0^TR_A(t)Q_0.
\]
The Woodbury identity, including singular $C$, gives
\begin{align}
 H_t&=(I+M_tC)^{-1}=I-Q_0^TR_B(t)Q_0C,
             &\|H_t\|&\le1+a/t,\label{eq:Hcore}\\
 K_t&=CH_t=(I+CM_t)^{-1}C=K_t^T,
             &D_t&:=R_A(t)-R_B(t)=V_tK_tV_t^T.
                                                        \label{eq:core}
\end{align}
$H_t$ need not be symmetric. The invertibility follows from the
determinant lemma and positive shifted endpoints. No inverse of $C$
occurs.

\begin{lemma}[Common inverse and core bounds]\label{lem:inverse}
One constant-distortion sparse subspace embedding of $X$ can be built
in $\widetilde O(Nm+m^3)$ conventional work. On one event of probability
at least $1-\delta$, it supports common symmetric linear approximations
$\widehat R_A(t)$ at all later shifts and right-hand sides, with
\begin{equation}\label{eq:inverse-error}
 \|\widehat R_A(t)-R_A(t)\|\le\nu_A/t.
\end{equation}
At each shift, setup is $O(m^3)$ and a block of $r$ right-hand sides
uses $O(Nmr\ell)$ work for $\ell=O(\log(1/\nu_A)+1)$ fixed polynomial
steps. If
\begin{equation}\label{eq:core-condition}
 \nu_A\le1,\qquad \nu_A(a/t)(1+a/t)\le\tfrac12,
\end{equation}
then substitution in \eqref{eq:core} is well defined and
\begin{equation}\label{eq:core-error}
 \|\widehat D_t-D_t\|\le
 \frac{6\nu_Aa(1+a/t)}{t^2}
       +\frac{2\nu_Aa^2(1+a/t)^2}{t^3}.
\end{equation}
The actual common symmetric operator
$\widehat R_B(t)=\widehat R_A(t)-\widehat D_t$ can therefore achieve
error $\nu_B/t$, $\nu_B\le1$, by the literal choice
\begin{equation}\label{eq:Bprecision}
 \nu_A\le\frac{\nu_B}{100(1+a/t)^4}.
\end{equation}
\end{lemma}

\begin{proof}
The embedding and fixed common inverse are \cite[Section~3]{Raw},
using the sparse embedding of \cite[Theorem~9]{NN}. In detail a
sketched Gram $G_0$ satisfies
$(3/4)G_0\preceq G=X^TX\preceq(5/4)G_0$. With $P_t=G_0+tI$,
the normalized error $P_t^{-1/2}(G-G_0)P_t^{-1/2}$ has norm at most
$1/4$. An even truncated Neumann series yields the stated common
symmetric physical inverse, uniformly in all right-hand sides. Square
roots of $P_t$ are analysis only; Cholesky and preconditioned polynomial
steps implement it. Its norm error decreases as $4^{-\ell}/t$.

For the core, $\|\widehat M_t-M_t\|\le\nu_A/t$. Inverse
perturbation using \eqref{eq:Hcore} gives
$\|\widehat H_t\|\le2(1+a/t)$ and
$\|\widehat K_t-K_t\|\le2\nu_Aa^2(1+a/t)^2/t$.
Expanding the two $V_t$ factors and then the core proves
\eqref{eq:core-error}; this is also \cite[Lemma~2.1]{Raw}.
The operator $\widehat D_t$ can be applied to any new block $W$ as
$\widehat V_t\widehat K_t(\widehat V_t^TW)$, using only thin
operations. Thus the bound for $B$ is a common-operator guarantee.
Writing $\Gamma=a/t$, \eqref{eq:core-error} implies
\[
 t\|\widehat R_B-R_B\|
 \le\nu_A(1+6\Gamma+8\Gamma^2+4\Gamma^3+2\Gamma^4)
 \le\nu_B
\]
under \eqref{eq:Bprecision}, which also ensures
\eqref{eq:core-condition}.
\end{proof}

For $c_\alpha=\sin(\pi\alpha)/\pi$, spectral calculus gives
\begin{equation}\label{eq:frac-integral}
 B^\alpha-A^\alpha=c_\alpha\int_0^\infty t^\alpha D_t\,dt,
 \qquad \|D_t\|\le\min(1/t,a/t^2).
\end{equation}
Both zero eigenvalues and signed $E$ are allowed. The norm bound and
the integral follow directly from positive endpoint resolvents and
their difference identity.

We use the positive quadrature in \cite[Lemma~5.1]{Raw} in its actual
scalar realization. In $t=ae^y$, the normalized logarithmic integrand
is analytic on $|\Im y|<\pi$ with envelope a fixed sector constant
times $e^{-g(y)}$, where
\begin{equation}\label{eq:g}
 g(y)=\alpha(-y)_++(1-\alpha)y_+.
\end{equation}
For $L\ge1$ the construction has $O_\alpha(L^2)$ positive actual
nodes $t_j=as_j$ and nonnegative logarithmic coefficients $a_j$,
with normalized error $O_\alpha(e^{-L})$ and
\begin{equation}\label{eq:quadrature-properties}
 |\log s_j|\le C_\alpha(L+1),\qquad
 \sum_j a_j\le C_\alpha(L+1).
\end{equation}
Its window length is $\ell=L^{3/4}+O_\alpha(1)$, enlargement
$b=O(\sqrt L)$, and local density
$[L-g(y)]_+/(2\pi d)$ with $d=\pi-L^{-1/4}$ up to certified
scalar-constant errors. Its actual log nodes differ from the ideal
lattice by $e^{-C_\alpha L}$. The scalar construction costs
$O_\alpha(L^3)$, using positive polynomial-square Gaussian windows
and Taylor/repeated-squaring exponential approximations. It uses no
input-dependent exponential, error function or best-pole oracle.
Choosing
\begin{equation}\label{eq:Lchoose}
 L=\log(1/\varepsilon)+O_\alpha(\log(2r))
\end{equation}
and sufficiently accurate common inverses gives the error
$\varepsilon\|E\|^\alpha$. These precise quadrature and inverse
inputs, rather than an uncharged approximation oracle, are the
imports from the earlier paper.

\section{Constructive sharp width at bounded endpoint scale}

We prove Theorem~\ref{thm:bounded}. Set $T=2Ra$, so
$\|A\|\le T/2$. Split the actual positive quadrature nodes into
low shifts $t<T$ and high shifts $t\ge T$.

\begin{lemma}[The low-node prefix]\label{lem:prefix}
The number of low actual nodes is bounded by
\begin{equation}\label{eq:low-count}
 q_{\rm low}\le\frac{L^2}{4\pi^2\alpha}
                             +O_{\alpha,R}(L^{7/4}).
\end{equation}
\end{lemma}
\begin{proof}
Their log coordinates lie below $\log(2R)$ up to the certified tiny
node perturbation. Every retained lattice point lies within $b$ of
its window. A low point therefore belongs to a window whose left
endpoint is at most $\log(2R)+b+1$. Count all these windows completely.
The extra prefix-boundary length is $O(\ell+b+1)$; its density is
$O(L+\ell)$. The same Riemann-sum and enlargement estimates as in
the full quadrature count give
\[
 q_{\rm low}\le\frac1{2\pi d}
     \int_{-\infty}^{\log(2R)}[L-g(y)]_+\,dy
       +O_{\alpha,R}(L\ell+L^2b/\ell).
\]
The negative half-line contributes exactly $L^2/(2\alpha)$, and
the fixed positive interval contributes $O_R(L+1)$. Since
$\ell=O(L^{3/4})$, $b=O(\sqrt L)$ and $d\uparrow\pi$ at rate
$L^{-1/4}$, the losses are $O_{\alpha,R}(L^{7/4})$.
\end{proof}

Each low node is computed by the previous common inverse and contributes
at most $r$ signed columns. At high shifts, instead use a single
polynomial block. Choose a sufficiently large even integer
\begin{equation}\label{eq:k-high}
 k_h\ge C_\alpha\bigl(L+\log(L+2)+1\bigr),\qquad k_h=O_\alpha(L).
\end{equation}
This is literal in the arithmetic model: for $L\ge1$,
$\log(L+2)\le L+1$, so an even integer above a larger fixed multiple
of $L+1$ suffices and is found by comparisons. Form
\begin{equation}\label{eq:high-W}
 W=[Q_0,(A/T)Q_0,\ldots,(A/T)^{k_h-1}Q_0].
\end{equation}
For $t\ge T$ set
\begin{equation}\label{eq:high-inverse}
 \widetilde R_A(t)=t^{-1}\sum_{j=0}^{k_h-1}(-A/t)^j.
\end{equation}
For even $k_h$ and $0\le x\le1/2$,
$\sum_{j=0}^{k_h-1}(-x)^j=(1-x^{k_h})/(1+x)$. Consequently
\begin{equation}\label{eq:high-inverse-bound}
 0\preceq\widetilde R_A(t)\preceq R_A(t),\qquad
 \|\widetilde R_A(t)-R_A(t)\|\le2^{-k_h}/t.
\end{equation}
In particular all high-node vectors have the exact common representation
\begin{equation}\label{eq:high-coeff}
 \widetilde V_t=W\cB_t,\qquad
 \cB_t=t^{-1}[I,-(T/t)I,\ldots,(-T/t)^{k_h-1}I]^T.
\end{equation}
Compute the small cores
$\widetilde M_t=Q_0^TW\cB_t$ and
$\widetilde K_t=C(I+\widetilde M_tC)^{-1}$. Since
$a/t\le1/(2R)<1/2$, Lemma~\ref{lem:inverse} gives
\begin{equation}\label{eq:high-node-error}
 \|\widetilde V_t\widetilde K_t\widetilde V_t^T-D_t\|
                                      \le C2^{-k_h}a/t^2.
\end{equation}
The physical quadrature weights $w_t$ are nonnegative. Sum all high
nodes into the one symmetric core
\begin{equation}\label{eq:high-core}
 J_h=\sum_{t\text{ high}}w_t\cB_t\widetilde K_t\cB_t^T.
\end{equation}
The corresponding correction $WJ_hW^T$ has width at most $rk_h$,
regardless of the number of high nodes. No basis of the full $X$
range is constructed.

Multiplying \eqref{eq:high-node-error} by $t^{\alpha+1}$ and using
$t\ge T\ge a$ bounds the normalized logarithmic node error by
$C_\alpha2^{-k_h}a^\alpha$. The coefficient mass in
\eqref{eq:quadrature-properties} is $O_\alpha(L+1)$, so
\eqref{eq:k-high} makes their total error fit in $a^\alpha e^{-L}$.
A fixed adjustment of the budget in \eqref{eq:Lchoose} absorbs this
along with the unchanged quadrature and low-solve errors. The width is
therefore at most
\[
 rq_{\rm low}+rk_h\le\frac{rL^2}{4\pi^2\alpha}
                               +O_{\alpha,R,r}(L^{7/4}),
\]
proving \eqref{eq:sharpwidth}.

Forming $W$ costs $O(Nmrk_h)$ through $X(X^T\cdot)$ actions, and
$Q_0^TW$ costs $O(Nr^2k_h)$. Even naive assembly and congruence
factorization of the $rk_h$ core cost at most $m^3$ times a fixed
polynomial in $L$. Multiplication of its final factor through $W$
costs $O(Nr^2k_h^2)$. Both are within the stated bound because
$r\le m\le N$. Adding the unchanged low nodes proves the work bound.

If $C\succeq0$, then $\widetilde M_t\succeq0$ and
$\widetilde K_t\succeq0$. If $C\preceq0$, write $C=-DD^T$
only in analysis. The exact shifted downdate gives
$I-D^TM_tD\succ0$, and \eqref{eq:high-inverse-bound} gives
$I-D^T\widetilde M_tD\succeq I-D^TM_tD$. Hence
$\widetilde K_t\preceq0$. The high sum thus has the claimed sign;
the low sum has the same sign property by \cite[Section~6]{Raw}.
Congruence elimination with scalar square roots returns all factors,
including singular cores. This completes the proof.

\begin{remark}
If only a computable upper bound $\Lambda$ is available, the same
algorithm may take $T=2\max(\Lambda,a)$, but its low-node prefix can
contain an additional $O(L\log(T/a))$ contribution. A Frobenius bound
that grows with dimension cannot be treated as a fixed $R$ in
\eqref{eq:sharpwidth}. The theorem instead uses its supplied promise
directly, setting $T/a=2R$.
\end{remark}

\section{Relative rank with singular endpoints}

In this section roots are analysis objects. Set
$P=A^{1/2}$, $Q_b=B^{1/2}$ and $d=\|\Delta\|$, with
$\Delta=Q_b-P$. Without any commutativity assumption,
\begin{equation}\label{eq:sylvester}
 E=P\Delta+\Delta Q_b.
\end{equation}
Therefore
\begin{equation}\label{eq:relative-lower}
 d\ge\frac{\|E\|}{2\sqrt\Lambda}\ge d_0.
\end{equation}
Also $d\le\sqrt{\|E\|}\le\sqrt a$, by Loewner--Heinz comparison,
so $\rho\ge1$. One may further use $d\le\sqrt\Lambda$, since
$0\preceq P,Q_b\preceq\sqrt\Lambda I$.

\begin{lemma}[Explicit relative ADI rank]\label{lem:adi-rank}
If $0<d_0\le d$, $\max(\|A\|,\|B\|)\le\Lambda$ and
$\rank E\le r$, an explicit low-rank ADI sum has width at most $kr$,
where
\[
 k=O\left(\log(C_*/\varepsilon)
   [\log(C_*/\varepsilon)+\log(\sqrt\Lambda/d_0)]\right),
\]
and error at most $\varepsilon d/2$. It allows singular endpoints.
\end{lemma}
\begin{proof}
Let $\lambda=\varepsilon d_0/8$, $\zeta=\varepsilon/8$ and
$M=\sqrt\Lambda$. Take the dyadic grid
$\tau_i=2^i\lambda$, $0\le i\le h$, where $h$ is least with
$2^h\lambda\ge M$. Repeat the entire grid $p$ times, where $p$ is
least with $3^{-p}\le\zeta$. There are $k=(h+1)p$ poles, each in
$[\lambda,2M]$, with the displayed bound. Doubling, multiplication by
three and comparisons compute the grid and repetition count.

Define $r_\tau(x)=(x-\tau)/(x+\tau)$ and
$\cR(x)=\prod_{j=1}^k r_{\tau_j}(x)$. For $x\ge0$,
$|\cR(x)|\le1$. For $\lambda\le x\le M$, one dyadic grid point
satisfies $\tau_i\le x\le2\tau_i$, giving $|r_{\tau_i}(x)|\le1/3$.
Thus $|\cR(x)|\le\zeta$ on $[\lambda,M]$.

Let $L_0$ be the spectral projection of $P$ on $[0,\lambda)$, and
$L_1$ that of $Q_b$; put $H_i=I-L_i$. The special identity
$\Delta=Q_b-P$ gives
\begin{equation}\label{eq:lowlow}
 \|L_0\Delta L_1\|\le2\lambda.
\end{equation}
Indeed the two terms $L_0Q_bL_1$ and $L_0PL_1$ each have norm at
most $\lambda$, even though the projections need not commute.
Decomposing $\Delta=L_0\Delta L_1+H_0\Delta+L_0\Delta H_1$
therefore gives
\begin{equation}\label{eq:adi-error}
 \|\cR(P)\Delta\cR(Q_b)\|
 \le2\lambda+2\zeta d\le\varepsilon d/2.
\end{equation}

For one pole, multiplying on the left by $P+\tau I$ and on the
right by $Q_b+\tau I$ proves the classical ADI identity
\begin{equation}\label{eq:adi-one}
 \Delta-r_\tau(P)\Delta r_\tau(Q_b)
       =2\tau(P+\tau I)^{-1}E(Q_b+\tau I)^{-1}.
\end{equation}
Set $U_0=V_0=Q_0$ and define
\begin{align}
 L_j&=(P+\tau_jI)^{-1}U_{j-1},&
 R_j&=(Q_b+\tau_jI)^{-1}V_{j-1},\label{eq:adi-factors}\\
 U_j&=U_{j-1}-2\tau_jL_j,&
 V_j&=V_{j-1}-2\tau_jR_j.\nonumber
\end{align}
The width-$kr$ sum
\begin{equation}\label{eq:adi-sum}
 F_k=\sum_{j=1}^k2\tau_jL_jCR_j^T
\end{equation}
satisfies $\Delta-F_k=\cR(P)\Delta\cR(Q_b)$ by telescoping
\eqref{eq:adi-one}. This proves the lemma. Its proof needs no
separation of the two root spectra at zero.
\end{proof}

The root Sylvester identity and rational-displacement/ADI mechanism
are established antecedents \cite{SDA,BT}. The estimate
\eqref{eq:lowlow} supplies the cutoff qualification when the positive
spectra touch at zero. The next section removes the apparently supplied
root actions in \eqref{eq:adi-factors}.

\section{Executing root ADI through raw shifted solves}

\begin{lemma}[Common inverse-root action]\label{lem:root-action}
For any PSD $S$, $\tau>0$ and $0<\eta<1$, positive quadrature over
ordinary shifted $S$ inverses constructs a common symmetric operator
$\widehat F_{S,\tau}$ satisfying
\begin{equation}\label{eq:root-action-error}
 \|\widehat F_{S,\tau}-(\sqrt S+\tau I)^{-1}\|\le\eta/\tau.
\end{equation}
It uses $O(\log^2(C_*/\eta))$ positive shifts with
$|\log(t/\tau^2)|\le C\log(C_*/\eta)$. Scalar work is a fixed
polynomial in that logarithm, using field operations and square roots.
\end{lemma}
\begin{proof}
Scalar integration after $t=u^2$, followed by spectral calculus, gives
\begin{equation}\label{eq:root-stieltjes}
 (\sqrt S+\tau I)^{-1}
  =\frac1\pi\int_0^\infty
           \frac{\sqrt t}{t+\tau^2}(S+tI)^{-1}\,dt.
\end{equation}
The identity remains valid at a zero eigenvalue. Substituting
$t=\tau^2e^y$, the logarithmic integrand has norm at most
\begin{equation}\label{eq:root-envelope}
 \frac1{\pi\tau}\frac{e^{y/2}}{1+e^y}
                         \le\frac1{\pi\tau}e^{-|y|/2}.
\end{equation}
It is analytic on $|\Im y|<\pi$; both denominators obey uniform
sector estimates. On the fixed half-strip $|\Im y|\le\pi/2$ the
boundary integral is $O(1/\tau)$. The infinite trapezoid error is at
most a constant times $\tau^{-1}e^{-\pi^2/h}$ for mesh $h$; truncating
at $|y|\le Y$ contributes $O(\tau^{-1}e^{-Y/2})$. Take
$h$ comparable to $1/\log(C_*/\eta)$ and
$Y=C\log(C_*/\eta)$. This proves the node and ideal error bounds.

The physical weights are
$h\,t\sqrt t/[\pi(t+\tau^2)]$. The same certified scalar
construction as in \cite[Lemma~5.1]{Raw} realizes them with positive
actual nodes and weights: approximate $\pi$ by its elementary series,
and $e^y$ by Taylor approximation at a small argument followed by
repeated squaring. Normalized real-line derivative bounds from
\eqref{eq:root-envelope} make the needed node precision a fixed
polynomial in $\eta$. This costs $O(\log^3(C_*/\eta))$ scalar
operations after a fixed error allocation, with no transcendental or
polynomial-root oracle.

If each physical inverse has error at most $\nu/t$, the total added
error is at most $C\nu/\tau$, because the positive quadrature of the
norm envelope \eqref{eq:root-envelope} has uniformly bounded mass.
Choose $\nu$ a sufficiently small fixed multiple of $\eta$ and split
the remaining scalar and truncation budgets. Lemma~\ref{lem:inverse}
provides these actual common inverse operators for either $S=A$ or
$S=B$, using \eqref{eq:Bprecision} for the latter.
\end{proof}

Execute every action in \eqref{eq:adi-factors} with
Lemma~\ref{lem:root-action} at the fixed tolerance
\begin{equation}\label{eq:eta}
 \eta=\frac{\varepsilon^2}{10^6\rho(k+1)^2}.
\end{equation}
All accuracy and loop caps can be selected before the adaptive products.
Since $\lambda\le\tau_j\le2M$, \eqref{eq:relative-scales} gives
\begin{equation}\label{eq:tau-range}
 \lambda^2/a=\varepsilon^2/(64\rho),\qquad
 M^2/a=\Lambda/a=\rho/(4r).
\end{equation}
Every physical shift in a root action consequently satisfies
\begin{equation}\label{eq:shift-range}
 |\log(t/a)|\le C\bigl(L+\log\rho+\log(C_*/\eta)\bigr).
\end{equation}
The precision \eqref{eq:Bprecision} thus adds only these logarithms to
the fixed Neumann degree. No small endpoint eigenvalue enters.

We prove error propagation explicitly. Each exact Cayley operator
$I-2\tau(\sqrt S+\tau I)^{-1}$ is a contraction; its computed operator
differs by at most $2\eta$. Ordered-product comparison, which does
not require computed operators to commute, gives for $j\le k$
\begin{align}
 \|\widehat U_j\|,\|\widehat V_j\|&\le(1+2\eta)^j,\label{eq:Uhatnorm}\\
 \|\widehat U_j-U_j\|,\|\widehat V_j-V_j\|
     &\le(1+2\eta)^j-1\le4k\eta.\label{eq:Uhaterr}
\end{align}
The chosen \eqref{eq:eta} ensures $2k\eta<2\cdot10^{-6}$ and
\begin{equation}\label{eq:strongfactor}
 \|\widehat L_j\|,\|\widehat R_j\|
 \le\frac{(1+\eta)(1+2\eta)^k}{\tau_j}<\frac2{\tau_j}.
\end{equation}
Combining \eqref{eq:root-action-error} with \eqref{eq:Uhaterr} yields
\begin{equation}\label{eq:factor-error}
 \|\widehat L_j-L_j\|,\|\widehat R_j-R_j\|
                         \le\frac{6(k+1)\eta}{\tau_j}.
\end{equation}
Expanding first in the left factor and then the right, and using
$\|C\|\le a$, gives the literal term bound
\[
 \|2\tau_j(\widehat L_jC\widehat R_j^T-L_jCR_j^T)\|
                 \le\frac{36a(k+1)\eta}{\tau_j}.
\]
Summing and inserting \eqref{eq:eta}, \eqref{eq:relative-scales} and
$\lambda=\varepsilon d_0/8$ gives
\begin{equation}\label{eq:ADI-inexact-error}
 \|\widehat F_k-F_k\|
 \le\frac{72a(k+1)^2\eta}{\lambda}
 =\frac{576}{10^6}\varepsilon d_0<\varepsilon d_0/2.
\end{equation}
Together with Lemma~\ref{lem:adi-rank}, this proves the full relative
error in Theorem~\ref{thm:relative}.

The Stieltjes stages act on $N$-by-$r$ blocks; they do not add final
factor columns. Thus $\widehat F_k$ has width at most $kr$. If
$\widehat F_k=P_fQ_f^T$, the symmetric signed representation is
\begin{equation}\label{eq:sym}
 \frac{\widehat F_k+\widehat F_k^T}{2}
 =\left(\frac{P_f+Q_f}{2}\right)
       \left(\frac{P_f+Q_f}{2}\right)^T
  -\left(\frac{P_f-Q_f}{2}\right)
       \left(\frac{P_f-Q_f}{2}\right)^T.
\end{equation}
It has no larger error to the symmetric target $\Delta$.

We finally charge the work. With $H$ as in the theorem,
$k=O(H^2)$ and $\log(C_*/\eta)=O(H)$. Each root action uses
$O(H^2)$ physical shifts; each physical inverse uses $O(H)$ raw
polynomial steps on a thin block and $O(m^3)$ setup. In addition to
the intrinsic reduction and one $\widetilde O(Nm+m^3)$ embedding,
an explicit overbound is
\begin{equation}\label{eq:cost-count}
 O(NmrH^5+m^3H^4+H^5).
\end{equation}
The scalar term is absorbed by $NmrH^5$. Costs $Nr^2$, $m^2r$ and
$r^3$ per stage are absorbed since $r\le m\le N$. Common inverse
operators make all estimates uniform in the adaptively generated
right-hand sides on the one embedding event; no failure union over
ADI stages is needed. On a bad event, failed small-core inverses or
required sign tests return zero, while all loops retain their caps.
This proves \eqref{eq:relative-work} with an every-outcome cost bound.
Internal Stieltjes nodes may be streamed; output storage is $O(Nkr)$
apart from the embedding/core workspace. No root, eigendecomposition
or Sylvester solve was assumed free.

\section{Deterministic preprocessing and scope}

We prove Theorem~\ref{thm:diagonal}. Compute $D$ in $O(Nm)$ and use
$G_0=D$ in Lemma~\ref{lem:inverse}. The promise
\eqref{eq:diagonal-promise} supplies the same constant Gram sandwich.
The shifted coarse matrix $D+tI$ is diagonal even if $D$ has zeros.
Every true Gram action is still $X^T(XW)$ on a thin block, costing
$O(Nmr)$. The same even Neumann proof therefore constructs all
common inverse operators without either an embedding or a dense
$m$-dimensional factorization. The only cubic cores are of update
size, and $r^3\le Nmr$. The count \eqref{eq:cost-count} becomes
$O(NmrH^5)$.

The bounded-spectrum construction uses the same diagonal coarse
inverse at low shifts and its unchanged common polynomial block at
high shifts. Its high core of size $rk_h$ also costs only $Nmr$
times accuracy logarithms because $r^3\le Nmr$. This proves the
second claim in Theorem~\ref{thm:diagonal}.

The promise controls correlations after column normalization, not
column magnitudes. It allows arbitrary positive diagonal rescaling
of otherwise well-preconditioned columns, and hence unbounded condition
numbers. It does not assert that a diagonal proxy can be produced for
an arbitrary raw factor at this cost. The cubic-free statement remains
a theorem for that explicit class.

\subsection*{Scale dependence and comparison}

The $\rho$ dependence in the relative theorem cannot simply disappear.
For any $n$, let $P=\diag(R^i)_{i=0}^{n-1}$, $R\ge289$, and
$u_i=\sqrt{2R^i}$. Set $A=P^2$, $E=\tau uu^T$ with $\tau=1/10$,
and $B=A+E$. The solution of $PT+TP=uu^T$ has entries
\[
 T_{ij}=\frac{2\sqrt{R^iR^j}}{R^i+R^j},\qquad
 T_{ii}=1,\qquad
 \sum_{j\ne i}|T_{ij}|\le\frac4{\sqrt R-1}\le\tfrac14.
\]
Thus $(3/4)I\preceq T\preceq(5/4)I$. The exact root correction
$\Delta$ satisfies $P\Delta+\Delta P+\Delta^2=\tau uu^T$, so
\[
 \tau T-\Delta=\int_0^\infty e^{-sP}\Delta^2e^{-sP}\,ds\succeq0.
\]
Since $\lambda_{\min}(P)=1$, one obtains
$\|\Delta\|\le5\tau/4$ and
$\|\tau T-\Delta\|\le25\tau^2/32$. Consequently
$\lambda_{\min}(\Delta)>\|\Delta\|/2$. Every rank-$<n$ approximation
has relative error greater than $1/2$. Here
$a=\Theta(\tau R^{n-1})$, $\Lambda=\Theta(R^{2(n-1)})$, and
$\log\rho=\Theta(n)$ at fixed $R,\tau$. The linear logarithmic scale
dependence in Theorem~\ref{thm:relative} has the correct order for
this family at fixed accuracy. This is neither an arithmetic lower
bound nor joint optimality in $\varepsilon,\rho$.

Shmueli--Drineas--Avron \cite{SDA} already use the root Sylvester
identity and relative singular-value decay, with an SPD separation
parameter involving $\lambda_{\min}(A)$. The singular-endpoint relative
rank order here also follows by a short regularization argument from
that published theorem: add $\gamma I$ to both endpoints with
$\gamma=(\varepsilon d_0/16)^2$. The root correction changes by at
most $2\sqrt\gamma\le\varepsilon d_0/8$, while the SPD estimate gives
the same $O(rL[L+\log\rho])$ rank order. Thus the low--low cutoff
proof is a direct alternative, rather than a new rank-order principle.
Regularized classical ADI also uses the same raw Stieltjes backend,
with physical shifts $t$ replaced by $t+\gamma$. The fully charged
implementation is a separate synthesis and complexity question.
Beckermann--Townsend
\cite{BT} establishes the rational-displacement rank mechanism and
its relation to ADI. Beckermann et al.\ \cite{BCKS} establishes rational
and polynomial matrix-function update methods. The common high-shift
span used here belongs to that classical framework. Classical positive
quadrature, root-exponential compression and sparse preconditioning
also predate these results \cite{BKS,NN,FTU}.

The precise consequences proved here are constructive sharp leading
width in the bounded endpoint class, and fully charged relative
principal-root corrections with a computable scale cutoff and singular
endpoints. A bounded primary-source search did not locate identical
statements; this is not a novelty certificate or an all-method
state-of-the-art comparison. No PSD-output claim for ADI, numerical
bit-complexity theorem, or unrestricted quadratic algorithm is implied.

\begin{thebibliography}{99}
\bibitem{Raw}
Working research manuscript,
\emph{Sharp-rank algorithms for low-rank updates of fractional matrix
functions}, 2026. Companion source:
\texttt{sharp-matrix-function-updates.tex}.

\bibitem{Bounded}
Working research note,
\emph{Sharp fractional compression with a bounded endpoint-to-update
ratio}, 2026. Companion source:
\texttt{bounded-spectrum-fractional-counts.txt}.

\bibitem{SDA}
S.~Shmueli, P.~Drineas and H.~Avron,
Low-rank updates of matrix square roots,
\emph{Numer. Linear Algebra Appl.} \textbf{31} (2024), e2528.
\href{https://doi.org/10.1002/nla.2528}{doi:10.1002/nla.2528}.
\href{https://arxiv.org/abs/2201.13156}{arXiv:2201.13156}.

\bibitem{BT}
B.~Beckermann and A.~Townsend,
On the singular values of matrices with displacement structure,
\emph{SIAM J. Matrix Anal. Appl.} \textbf{38} (2017), 1227--1248;
expanded as \emph{SIAM Rev.} \textbf{61} (2019), 319--344.
\href{https://doi.org/10.1137/16M1096426}{doi:10.1137/16M1096426}.

\bibitem{BCKS}
B.~Beckermann, A.~Cortinovis, D.~Kressner and M.~Schweitzer,
Low-rank updates of matrix functions II: Rational Krylov methods,
\emph{SIAM J. Numer. Anal.} \textbf{59} (2021), 1325--1347.
\href{https://doi.org/10.1137/20M1362553}{doi:10.1137/20M1362553}.

\bibitem{BKS}
M.~Sh.~Birman, L.~S.~Koplienko and M.~Z.~Solomyak,
Estimates of the spectrum of a difference of fractional powers of
selfadjoint operators, \emph{Izv. Vyssh. Uchebn. Zaved. Mat.}
(1975), no.~3, 3--10.
\href{https://www.mathnet.ru/eng/ivm6330}{Math-Net record}.

\bibitem{NN}
J.~Nelson and H.~L.~Nguy\^en,
OSNAP: Faster numerical linear algebra algorithms via sparser subspace
embeddings, \emph{Proc. FOCS} (2013), 117--126.
\href{https://arxiv.org/abs/1211.1002}{arXiv:1211.1002}.

\bibitem{FTU}
Z.~Frangella, J.~A.~Tropp and M.~Udell,
Randomized Nystr\"om preconditioning,
\emph{SIAM J. Matrix Anal. Appl.} \textbf{44} (2023), 718--752.
\href{https://doi.org/10.1137/21M1466244}{doi:10.1137/21M1466244}.
\end{thebibliography}
\end{document}
